% chapters/theory/06.tex

\chapter{Reasoning Models and Verifiable Reinforcement Learning}
\label{chap:reasoning-models}

In the early wave of instruction-tuned assistants, many users discovered a sharp boundary: models could write fluent explanations and generate plausible code, yet still fail on problems that require careful multi-step reasoning. The last year has seen a rapid shift toward “reasoning models,” systems that explicitly allocate computation to produce intermediate reasoning before committing to an answer. This chapter develops a practical understanding of what reasoning means in the LLM context, how reasoning is evaluated, and why reinforcement learning becomes especially powerful when the task admits a verifiable reward.

\section{What we mean by reasoning}
There is no universally agreed definition of reasoning in language models, but there is a useful operational one. A question is “reasoning-heavy” when it cannot be answered reliably by simple recall or pattern completion, and instead requires a sequence of intermediate steps that constrain the final answer. In human terms, it resembles an exam question where the solution emerges from a chain of deductions rather than a single remembered fact.

A knowledge question might ask for a course code. A reasoning question might ask for the age of a bear given its birth year and the current year. That example is trivial, but it captures the structure: there is a latent calculation separating the prompt from the answer. In more realistic cases—contest math, algorithmic programming, or multi-hop logic—mistakes can compound. One incorrect step early in the chain can poison everything that follows.

\section{Reasoning as deliberate intermediate text}
A simple conceptual precursor to reasoning models is chain-of-thought prompting: ask the model to “think step by step.” The core idea is that the model should produce intermediate reasoning tokens before outputting a final answer. Even without changing weights, this often improves performance because it gives the model more opportunities to condition on partially structured progress, rather than jumping directly to a final token sequence.

Reasoning models systematize this behavior. They are trained so that their default response structure is not merely an answer, but a reasoning segment followed by an answer segment. In many user interfaces, the full reasoning chain is hidden and only a short summary is displayed. Two practical motivations are common. First, long chains can be hard to read and are often not necessary for the user to act on the final answer. Second, reasoning tokens are still tokens; in API settings they count toward output cost. A system that reasons well but wastes tokens is expensive to serve.

\section{Benchmarks for reasoning}
Reasoning is not evaluated by “vibes.” It is measured on tasks where correctness can be determined by rules or tests. This is crucial: it gives us objective feedback and, later, verifiable reinforcement signals.

For coding, a common benchmark structure is straightforward. The model receives a problem statement and must produce a program that passes a hidden test suite. HumanEval is a canonical dataset of this form, with a moderate number of human-written problems. Codeforces-style benchmarks extend this to competitive programming. SWE-bench goes further by grounding tasks in real software engineering: the model is given a repository and a bug report, and it must produce a patch. Correctness is determined by running tests.

For mathematics, benchmarks such as AIME provide difficult contest-style questions with constrained answer formats. GSM8K provides grade-school math word problems. In both cases, the final answer is typically easy to parse if the prompt enforces a format, and the answer can be checked against ground truth.

\section{pass@k: correctness under multiple attempts}
Reasoning models are often sampled multiple times. The reason is pragmatic: for tasks with a verifiable checker (tests for code, ground truth for math), we can afford to try several solutions and accept success if any attempt is correct. This motivates a metric called pass@k.

Pass@k is the probability that at least one of \(k\) attempts succeeds. If a model produces \(k\) candidate solutions, pass@k measures the chance that at least one is correct. This is meaningful when generation is stochastic and correctness can be verified automatically.

A naïve estimate would be to sample \(k\) times and count success. But that estimate is noisy. A standard approach instead samples \(n\) total solutions for a prompt, observes that \(c\) of them are correct, and then estimates what the probability would be that at least one of \(k\) randomly selected attempts is correct.

To derive this estimator, it is easier to compute the complementary event: none of the \(k\) selected attempts are correct. If there are \(n-c\) incorrect solutions among the \(n\) total, then sampling without replacement gives
\[
P(\text{all } k \text{ are incorrect}) \;=\; \frac{n-c}{n}\cdot \frac{n-c-1}{n-1}\cdots \frac{n-c-k+1}{n-k+1}.
\]
This product can be expressed using combinations:
\[
P(\text{all } k \text{ incorrect}) \;=\; \frac{\binom{n-c}{k}}{\binom{n}{k}}.
\]
Therefore the standard estimator is
\[
\mathrm{pass@k} \;=\; 1 - \frac{\binom{n-c}{k}}{\binom{n}{k}}.
\]
The special case \(k=1\) reduces to the empirical accuracy \(c/n\), which matches intuition.

Sampling temperature matters here. If temperature is near zero, repeated samples are nearly identical and pass@k barely increases with \(k\). If temperature is moderate, diversity increases and pass@k can improve substantially. If temperature is too high, quality deteriorates and pass@k can worsen. This is why benchmark reports almost always specify temperature.

\section{Why reinforcement learning is a natural fit}
If the goal is to train a model to produce long, correct reasoning chains, one could attempt supervised learning: collect prompts paired with “ideal” reasoning traces. In practice this is hard for three reasons. First, producing high-quality reasoning traces is expensive and time-consuming. Second, the best internal reasoning style for a model need not match the style humans naturally write. Third, in the most important reasoning domains, we often have something better than a human-written chain: we have a checker.

For math tasks, we can verify the final answer against ground truth. For coding tasks, we can run tests. This gives a verifiable reward: a completion is correct or incorrect. When such a reward exists, reinforcement learning becomes attractive. We can use the model itself to explore possible solution paths, and then use verification to provide a learning signal without needing to label every intermediate step.

In practice, reasoning training often uses a reward that combines two ingredients. One ingredient enforces the presence of a reasoning segment, for example by checking that the model produced specific “think” delimiters. The other ingredient checks final correctness. The first ingredient ensures the behavior we want—deliberate reasoning—and the second ensures that the reasoning leads somewhere useful.

This setup produces a striking empirical pattern: as RL training progresses, benchmark performance improves. In other words, models can discover and refine reasoning behaviors simply by being rewarded for correct answers and for producing reasoning traces.

\section{Budgeting and “overthinking”}
Once a model is trained to reason, a new problem appears: not every prompt deserves long reasoning. Some tasks are trivial, and long reasoning is wasted latency and cost. This motivates dynamic control of reasoning length. One approach is to classify prompts into “easy” and “hard” categories and allocate different reasoning budgets. Another is budget forcing: inject tokens that encourage the model to continue thinking, or inject a signal that forces it to stop and answer. There is also active research on moving reasoning into continuous latent representations rather than explicit tokens, compressing the “thinking” process without consuming as many output tokens.

\section{GRPO: Group Relative Policy Optimization}
Preference tuning introduced PPO as a widely used policy optimization algorithm, but reasoning training in 2024--2025 popularized a different family of methods. GRPO—Group Relative Policy Optimization—shares the same high-level intent as PPO: improve policy behavior while controlling drift. The key difference lies in how the advantage signal is computed.

PPO typically learns a value function. The value function predicts the expected return from a partial state (the prompt plus partial generation). The advantage is then computed as reward minus baseline, often through generalized advantage estimation. This introduces an extra model head to train and maintain, and it adds complexity.

GRPO removes this component. Instead of predicting a baseline with a value function, GRPO constructs a baseline from the model’s own samples. For a given prompt, it generates multiple completions—call the group size \(g\). Each completion receives a reward. The advantage for a completion is then computed relative to the group, typically by subtracting the group mean and normalizing by the group standard deviation:
\[
A_i \;=\; \frac{r_i - \mu(r)}{\sigma(r)}.
\]
This has a natural interpretation. If a problem is easy, many completions will score well, and the group mean will be high; a single correct completion is not as informative. If a problem is hard, most completions will fail, and a single correct completion stands out strongly. The group-relative advantage therefore adapts to difficulty without requiring a learned value function.

In reasoning settings, GRPO also benefits from verifiable rewards. Because correctness can be checked, there may be no need for a separate learned reward model at all. The training system can use a checker to compute reward directly, and a reference model to prevent policy drift through KL regularization.

\section{Length growth and length bias}
A consistent empirical observation in reasoning RL is that output length increases during training. Early on, this can be correlated with performance improvements: longer reasoning enables deeper exploration. But in later stages, performance can plateau while length continues to grow. This matters because longer outputs increase latency and cost.

Analyses of GRPO-style objectives identified a length-dependent normalization that changes token-level contributions depending on whether a token appears in a short or long output. The effect can unintentionally encourage the model to prefer longer incorrect outputs over shorter incorrect ones, because the penalty per token differs with length.

Several fixes have been proposed. One line of work normalizes token contributions so they are comparable across outputs, removing the length-based incentive. Another removes the problematic normalization term entirely. These modifications tend to preserve the length of correct solutions while reducing unnecessary verbosity in incorrect solutions.

Additional refinements address edge cases in group normalization. When a prompt is extremely hard and most completions fail, the reward variance can behave poorly, making the standard deviation normalization unstable or uninformative. Some implementations adjust this normalization. Other refinements introduce asymmetric clipping bounds so that low-probability tokens can grow meaningfully without allowing high-probability tokens to collapse abruptly.

\section{A concrete training recipe: DeepSeek R1}
It is useful to see how these ideas assemble into an end-to-end pipeline. A prominent example is DeepSeek’s R1 line of work, which demonstrated strong reasoning capabilities and provided a relatively clear recipe.

The first stage, sometimes referred to as a “zero” variant, starts from a strong pretrained base model and applies reinforcement learning directly on reasoning tasks, using verifiable rewards. The reward is typically a combination of correctness and a formatting constraint that enforces a reasoning section followed by an answer section. This stage demonstrates that large gains in reasoning benchmarks can be achieved without supervised reasoning traces.

However, RL-only training can produce undesirable artifacts in the reasoning text itself—mixed languages, poor readability, inconsistent formatting. The full R1 pipeline addresses this by introducing a small amount of “cold start” supervised data. In this step, a subset of model-generated reasoning traces is corrected by humans to enforce clean formatting and language consistency. The model is first fine-tuned on this small curated set, which anchors its reasoning style.

After this anchoring, reinforcement learning is applied again on reasoning tasks, now with additional heuristics such as language consistency rewards. Next, the model undergoes another supervised stage that mixes reasoning and non-reasoning data. The reasoning portion can be generated via rejection sampling: the model produces multiple candidate solutions and only high-quality, correct, well-formatted solutions are kept. This improves coverage and stabilizes behavior. Finally, an additional alignment stage incorporates helpfulness and harmlessness objectives, ensuring that the model remains usable as an assistant while being safe, including within its internal reasoning traces.

\section{Distillation: transferring reasoning to smaller models}
A practical issue remains: many organizations cannot deploy extremely large reasoning models everywhere. One response is distillation. A large teacher reasoning model generates many examples with reasoning and final answers. A smaller student model is then trained via supervised learning to imitate these sequences. This is distillation in the sense of transferring behavior rather than matching probability distributions.

Empirically, for smaller models, distillation can be more effective than applying RL directly from scratch. RL requires exploration and a strong baseline policy; small models may struggle to discover high-quality reasoning trajectories without guidance. Distillation provides that guidance in the form of explicit sequences to imitate.

\section{Summary}
Reasoning models can be understood as models trained to produce intermediate reasoning before an answer, allocating more computation to hard tasks. Their performance is evaluated on benchmarks with verifiable outcomes, and pass@k formalizes the benefit of multiple sampled attempts when a checker exists. Reinforcement learning becomes especially powerful in this setting because correctness can be validated without human-written reasoning traces. GRPO is a central algorithm in this wave because it avoids value-function training by computing advantage relative to a group of sampled completions. Practical training reveals a tendency toward increasing output length, which motivates objective-function refinements that reduce length bias. A multi-stage recipe—small supervised anchoring, verifiable RL, mixture SFT, and final alignment—produces robust reasoning assistants, and distillation transfers reasoning behavior to smaller models efficiently.

\section{Exercises}
\paragraph{1. pass@k derivation.}
Re-derive \(\mathrm{pass@k} = 1 - \binom{n-c}{k}/\binom{n}{k}\) by computing the probability that all \(k\) selected samples are incorrect. State clearly what assumptions are being made.

\paragraph{2. Temperature and reliability.}
Explain why pass@k may not improve when temperature is near zero. Why can pass@k also degrade if temperature is too high?

\paragraph{3. Why verifiable rewards matter.}
Give two examples of tasks where correctness is verifiable and two examples where correctness is not easily verifiable. Explain how this changes the training approach.

\paragraph{4. GRPO vs PPO.}
Explain in prose the difference between GRPO and PPO in how they compute advantage. Why does GRPO not require a value function?

\paragraph{5. Length bias intuition.}
Describe how a length-dependent normalization term could encourage longer incorrect reasoning chains. What objective modification might reduce this incentive?
